% !TEX root =  Lecture14.tex


%%%%%%%%%%%%%%%%%%%%%%%%%%%%%%%%%%%%
% Recap/Agenda
%%%%%%%%%%%%%%%%%%%%%%%%%%%%%%%%%%%%
\begin{frame}
\frametitle{Module 3: Regression Inference}
    \begin{itemize}
        \item \hl{Previously:} inference about \emph{one} slope: the two roads, a bootstrap interval, LINE (Lecture 13).
        \item \hl{This time:} \hl{many predictors at once}. That slope test becomes one \emph{row} of an \texttt{lm()} table.
          \begin{itemize}
              \item reading the output: $\mathrm{df} = n-p-1$, every slope \emph{holding the others fixed};
              \item what \hl{multicollinearity} does to those SEs, $p$-values and intervals;
              \item when the conditions fail: \hl{transformations}, and reading a log model;
              \item \hl{interaction} terms.
          \end{itemize}
        \item \hl{Reading:} IMS Ch.\ 25.1--25.2, plus the non-linear and interaction supplements.
        \item \hl{Homework:} the assigned reading.
    \end{itemize}
\end{frame}


%%%%%%%%%%%%%%%%%%%%%%%%%%%%%%%%%%%%
% Learning objectives:
%%%%%%%%%%%%%%%%%%%%%%%%%%%%%%%%%%%%
\begin{frame}
    \frametitle{Lecture Objectives}
    \goals{%
    \begin{itemize}
    \item read a multiple regression coefficient table, and test $H_0:\beta_j=0$ against $t_{n-p-1}$;
    \item interpret each $\hat\beta_j$ \emph{holding the other predictors fixed}, and say why every row is conditional on the rest of the model;
    \item explain what \hl{collinear} predictors do to the inference, and why they leave prediction untouched;
    \item diagnose a failing model, \hl{transform} it, and interpret a $\log(y)$, $\log(x)$ or log--log coefficient;
    \item fit and read an \hl{interaction}: what $\hat\beta_3$ means when one slope depends on another predictor.
    \end{itemize}}
\end{frame}


%%%%%%%%%%%%%%%%%%%%%%%%%%%%%%%%%%%%
\begin{frame}
    \frametitle{Lecture Objectives}
    \objectives{%
    \begin{itemize}
    \item \textbf{(M3, LO3)} Interpreting Software Output for Multiple Regression
    \item \textbf{(M3, LO4)} Building and Interpreting Non-linear and Interaction Models
    \end{itemize}}
    \vspace{0.3em}
    {\small Also revisited, in the light of \emph{inference}:}
    \objectives{%
    \begin{itemize}
    \item \textbf{(M3, LO2)} Model Checking and Remedial Measures: the transformations of Lecture 13, now as model \emph{building}
    \item \textbf{(M1, LO3)} Using Regression Models: multicollinearity, from Lecture 7
    \end{itemize}}
\end{frame}
